% ===================================================================
% PHẦN NÀY DO: Nguyễn Hữu Gia Minh phụ trách
% Nội dung: HD Video Streaming với Frame Fragmentation (3 điểm)
% ===================================================================

\section{HD Video Streaming và Frame Fragmentation (3 điểm)}

\subsection{Vấn đề với UDP MTU}

Khi triển khai cơ bản (1 frame = 1 packet), hệ thống gặp vấn đề với HD video:

\begin{table}[H]
\centering
\begin{tabular}{|l|c|c|}
\hline
\textbf{Resolution} & \textbf{Average Frame Size} & \textbf{Vấn đề} \\
\hline
SD (640x480) & ~5-15 KB & OK với UDP MTU \\
HD (1280x720) & ~30-80 KB & Vượt quá MTU \\
Full HD (1920x1080) & ~100-200 KB & \textbf{CẦN FRAGMENTATION} \\
\hline
\end{tabular}
\caption{Frame Sizes theo Resolution}
\end{table}

\textbf{UDP MTU (Maximum Transmission Unit):}
\begin{itemize}
    \item Ethernet MTU: 1500 bytes
    \item IP header: 20 bytes
    \item UDP header: 8 bytes
    \item \textbf{UDP payload tối đa: 1472 bytes}
\end{itemize}

\textbf{Hậu quả khi frame size > MTU:}
\begin{itemize}
    \item Packets bị OS fragmentation ở IP layer
    \item Nếu 1 IP fragment bị mất → toàn bộ frame corrupted
    \item Client không thể detect và retransmit fragment cụ thể
    \item Video bị freeze hoặc artifacts nghiêm trọng
\end{itemize}

\subsection{Frame Fragmentation Strategy}

\subsubsection{Thiết kế giải pháp}

Chia mỗi frame thành nhiều RTP packets, mỗi packet ≤ MTU:

\begin{lstlisting}[language=Python, caption={Frame Fragmentation Constants}]
# Fragmentation parameters
MTU_SIZE = 1400         # Safe MTU size (< 1472)
HEADER_SIZE = 12        # RTP header size
CHUNK_SIZE = MTU_SIZE - HEADER_SIZE  # Payload per packet
\end{lstlisting}

\textbf{Fragment size calculation:}
\begin{itemize}
    \item MTU safe size: 1400 bytes (để dự phòng)
    \item RTP header: 12 bytes
    \item \textbf{Payload per fragment: 1388 bytes}
\end{itemize}

\subsubsection{Marker Bit Usage}

Marker bit trong RTP header được sử dụng để đánh dấu fragment cuối của frame:

\begin{table}[H]
\centering
\begin{tabular}{|c|l|l|}
\hline
\textbf{Marker bit} & \textbf{Ý nghĩa} & \textbf{Action} \\
\hline
0 & Fragment giữa & Lưu vào buffer \\
1 & Fragment cuối cùng & Reassemble frame \\
\hline
\end{tabular}
\caption{Marker Bit Semantics}
\end{table}

\textbf{Lưu ý:} Đây là điểm khác biệt chính so với triển khai cơ bản (nơi marker bit luôn = 0).

\subsection{Server-Side Fragmentation}

\subsubsection{makeRtp() - HD Version}

\begin{lstlisting}[language=Python, caption={Make RTP Packet with Fragmentation Support}]
def makeRtp(self, payload, frameNbr, marker=0):
    """
    Tao RTP packet voi ho tro fragmentation
    
    Args:
        payload: Fragment data (or complete frame)
        frameNbr: Sequence number cua packet
        marker: 1 neu la fragment cuoi, 0 neu la fragment giua
    
    Returns:
        Complete RTP packet (header + payload)
    """
    version = 2
    padding = 0
    extension = 0
    cc = 0
    # marker bit = 1 CHỈ KHI là fragment cuối cùng
    pt = 26  # MJPEG payload type
    seqnum = frameNbr
    ssrc = 0
    
    rtpPacket = RtpPacket()
    rtpPacket.encode(version, padding, extension, cc, 
                     seqnum, marker, pt, ssrc, payload)
    return rtpPacket.getPacket()
\end{lstlisting}

\textbf{Tham số mới:}
\begin{itemize}
    \item \texttt{marker}: Tham số để đánh dấu fragment cuối cùng
    \item Giá trị: 0 (fragment giữa), 1 (fragment cuối)
\end{itemize}

\subsubsection{sendRtp() - HD Version}

\begin{lstlisting}[language=Python, caption={Send RTP with Fragmentation}]
def sendRtp(self):
    """Main streaming loop voi frame fragmentation"""
    while True:
        frame_start = time.time()
        
        # Check for pause/stop
        if self.clientInfo['event'].isSet():
            print("Stopping RTP stream")
            break
        
        # Doc frame tu video file
        data = self.clientInfo['videoStream'].nextFrame()
        
        if not data:
            print("End of video")
            break
        
        frameNumber = self.clientInfo['videoStream'].frameNbr()
        
        try:
            address = self.clientInfo['rtspSocket'][1][0]
            port = int(self.clientInfo['rtpPort'])
            
            # FRAGMENTATION LOGIC
            frame_size = len(data)
            
            if frame_size <= CHUNK_SIZE:
                # Frame nho: gui 1 packet voi marker=1
                packet = self.makeRtp(data, frameNumber, marker=1)
                self.clientInfo['rtpSocket'].sendto(packet, 
                                                    (address, port))
                self.frames_sent += 1
                self.bytes_sent += frame_size
                
            else:
                # Frame lon: chia thanh nhieu fragments
                num_fragments = (frame_size + CHUNK_SIZE - 1) // CHUNK_SIZE
                print(f"Frame {frameNumber}: {frame_size}B -> "
                      f"{num_fragments} fragments")
                
                for i in range(num_fragments):
                    start = i * CHUNK_SIZE
                    end = min(start + CHUNK_SIZE, frame_size)
                    fragment = data[start:end]
                    
                    # Marker = 1 CHỈ cho fragment cuối cùng
                    marker = 1 if (i == num_fragments - 1) else 0
                    
                    # Sequence number tang dan cho moi fragment
                    seq = frameNumber + i
                    
                    packet = self.makeRtp(fragment, seq, marker=marker)
                    self.clientInfo['rtpSocket'].sendto(packet, 
                                                        (address, port))
                    
                    self.bytes_sent += len(fragment)
                
                self.frames_sent += 1
            
        except Exception as e:
            print(f"Send error: {e}")
            break
        
        # Frame rate control (25 FPS)
        frame_elapsed = time.time() - frame_start
        time.sleep(max(0, 0.04 - frame_elapsed))
\end{lstlisting}

\textbf{Fragmentation algorithm:}
\begin{enumerate}
    \item Kiểm tra frame size
    \item Nếu frame size ≤ CHUNK\_SIZE: gửi 1 packet với marker=1
    \item Nếu frame size > CHUNK\_SIZE:
    \begin{itemize}
        \item Tính số fragments: $\lceil \frac{\text{frame\_size}}{\text{CHUNK\_SIZE}} \rceil$
        \item Chia frame thành các fragments
        \item Gửi từng fragment với marker=0
        \item Fragment cuối cùng: marker=1
    \end{itemize}
    \item Sequence number tăng dần cho mỗi fragment
\end{enumerate}

\subsection{Client-Side Reassembly}

\subsubsection{listenRtp() - HD Version}

\begin{lstlisting}[language=Python, caption={Receive and Reassemble RTP Packets}]
def listenRtp(self):
    """Thread nhan RTP packets va reassemble frames"""
    self.frame_buffer = bytearray()  # Buffer cho fragments
    current_frame_seq = -1           # Sequence cua frame hien tai
    
    while True:
        try:
            if self.playEvent.isSet():
                break
            
            data = self.rtpSocket.recv(20480)  # Large buffer for fragments
            
            if data:
                rtpPacket = RtpPacket()
                rtpPacket.decode(data)
                
                currFrameNbr = rtpPacket.seqNum()
                marker = rtpPacket.marker()  # Doc marker bit
                
                # Neu la frame moi
                if currFrameNbr != current_frame_seq:
                    current_frame_seq = currFrameNbr
                    self.frame_buffer = bytearray()  # Reset buffer
                
                # Them fragment vao buffer
                self.frame_buffer += rtpPacket.getPayload()
                
                # NEU marker=1: day la fragment cuoi -> reassemble
                if marker == 1:
                    # Frame hoan thanh
                    complete_frame = bytes(self.frame_buffer)
                    
                    # Update metrics
                    self.updateMetrics(rtpPacket, len(complete_frame))
                    
                    # Hien thi frame
                    self.writeFrame(complete_frame)
                    
                    # Reset buffer cho frame tiep theo
                    self.frame_buffer = bytearray()
                    current_frame_seq = -1
                
        except socket.timeout:
            continue
        except Exception as e:
            if not self.playEvent.isSet():
                print(f"RTP receive error: {e}")
            break
\end{lstlisting}

\textbf{Reassembly algorithm:}
\begin{enumerate}
    \item Nhận RTP packet
    \item Decode và kiểm tra sequence number
    \item Nếu sequence khác frame hiện tại → reset buffer
    \item Append payload vào frame buffer
    \item Kiểm tra marker bit:
    \begin{itemize}
        \item Marker = 0: chờ fragment tiếp theo
        \item Marker = 1: frame hoàn thành → display
    \end{itemize}
    \item Reset buffer và chờ frame mới
\end{enumerate}

\subsubsection{Xử lý packet loss}

\begin{lstlisting}[language=Python, caption={Packet Loss Detection}]
def updateMetrics(self, rtpPacket, frameSize):
    """Cap nhat metrics va detect packet loss"""
    currSeq = rtpPacket.seqNum()
    
    # Detect packet loss
    if self.lastSeq != -1:
        expected = self.lastSeq + 1
        if currSeq != expected:
            lost = currSeq - expected
            self.packetsLost += lost
            print(f"Packet loss detected: {lost} packets")
    
    self.lastSeq = currSeq
    self.bytesRecv += frameSize
    self.framesRecv += 1
\end{lstlisting}

\textbf{Packet loss handling:}
\begin{itemize}
    \item So sánh sequence number hiện tại với expected
    \item Gap trong sequence = packet loss
    \item Đếm số packets lost
    \item Frame bị loss fragment → skip (chờ frame mới)
\end{itemize}

\subsection{Ví dụ Frame Fragmentation}

\textbf{Giả sử:} Frame Full HD với size = 150 KB

\begin{lstlisting}[caption={Fragmentation Example}]
Frame size: 153600 bytes
Chunk size: 1388 bytes
Number of fragments: ceil(153600 / 1388) = 111 fragments

Fragment 1:   Seq=1000, Marker=0, Payload[0:1388]
Fragment 2:   Seq=1001, Marker=0, Payload[1388:2776]
Fragment 3:   Seq=1002, Marker=0, Payload[2776:4164]
...
Fragment 110: Seq=1109, Marker=0, Payload[151472:152860]
Fragment 111: Seq=1110, Marker=1, Payload[152860:153600]  <- MARKER=1!
\end{lstlisting}

\textbf{Timeline client-side:}
\begin{enumerate}
    \item Nhận fragments 1-110: append vào buffer
    \item Nhận fragment 111 với marker=1
    \item Reassemble 111 fragments thành complete frame
    \item Display frame
    \item Reset buffer, chờ frame tiếp theo
\end{enumerate}

\subsection{Performance Metrics}

\subsubsection{Bandwidth Calculation}

\begin{lstlisting}[language=Python, caption={Calculate Statistics}]
def calculateStats(self):
    """Tinh bandwidth va thong ke"""
    if self.framesRecv > 0:
        elapsed = time.time() - self.startTime
        
        # Bandwidth (Mbps)
        bandwidth = (self.bytesRecv * 8) / (elapsed * 1e6)
        
        # Frame rate (FPS)
        fps = self.framesRecv / elapsed
        
        # Packet loss rate
        totalPackets = self.lastSeq + 1
        lossRate = (self.packetsLost / totalPackets) * 100
        
        print(f"Bandwidth: {bandwidth:.2f} Mbps")
        print(f"Frame rate: {fps:.2f} FPS")
        print(f"Packet loss: {lossRate:.2f}%")
\end{lstlisting}

\textbf{Metrics:}
\begin{itemize}
    \item \textbf{Bandwidth:} Bytes received × 8 / elapsed time
    \item \textbf{Frame rate:} Frames received / elapsed time
    \item \textbf{Loss rate:} Packets lost / total packets
\end{itemize}

\subsection{So sánh triển khai Basic vs HD}

\begin{table}[H]
\centering
\begin{tabular}{|l|c|c|}
\hline
\textbf{Feature} & \textbf{Basic (4đ)} & \textbf{HD (3đ)} \\
\hline
Frame size limit & < 1472 bytes & Unlimited \\
Fragments per frame & 1 & 1-200+ \\
Marker bit & Always 0 & 0 or 1 \\
Sequence number & Per frame & Per fragment \\
Client buffer & Direct display & Reassembly buffer \\
HD support & ❌ & ✅ \\
Packet loss detection & Basic & Advanced \\
\hline
\end{tabular}
\caption{Basic vs HD Implementation}
\end{table}

\subsection{Tóm tắt HD Video Streaming}

\subsubsection{Server-side}
\begin{itemize}
    \item Kiểm tra frame size
    \item Chia frame thành fragments nếu > MTU
    \item Set marker=1 cho fragment cuối
    \item Tăng sequence number cho mỗi fragment
    \item Gửi tất cả fragments qua UDP
\end{itemize}

\subsubsection{Client-side}
\begin{itemize}
    \item Nhận RTP packets liên tục
    \item Buffer fragments theo sequence number
    \item Detect marker=1 để reassemble frame
    \item Display frame hoàn chỉnh
    \item Detect packet loss qua sequence gaps
\end{itemize}

\textbf{Kết quả:} Hỗ trợ streaming Full HD (1920×1080) với frame size lên đến 200 KB mà không bị fragmentation ở IP layer.

\clearpage
